%%% FIELD proof-prop-unsubsampled-reduction
Formally specialize Definitions~\(\texttt{def:compatible-fiber}\) and \(\texttt{def:certificate}\) to \(m=1\). Let \((A,D)\in\mathcal S(B,Q;G)\). The first compatibility equation in Definition~\(\texttt{def:compatible-fiber}\) gives
\[
B=A^1=A.
\]
The second compatibility equation in that definition, together with \(A^0=I_d\), gives
\[
Q=\sum_{r=0}^{1-1}A^rD(A^r)^\top
  =A^0D(A^0)^\top
  =D.
\]
Thus every member of \(\mathcal S(B,Q;G)\) equals \((B,Q)\), so the fiber has at most one member.

Conversely, substitute \((A,D)=(B,Q)\) into Definition~\(\texttt{def:compatible-fiber}\). The two coarse equations hold identically because
\[
B^1=B
\qquad\text{and}\qquad
\sum_{r=0}^{0}B^rQ(B^r)^\top=B^0Q(B^0)^\top=Q.
\]
The remaining defining conditions hold if and only if
\[
B_{ij}=0\quad(i<j),
\qquad
B_{ij}=0\quad(i>j\text{ and }j\to i\notin G),
\]
\[
\max_i|B_{ii}|<1,
\qquad
Q_{ij}=0\quad(i\ne j),
\qquad
Q\succ0.
\]
These are precisely the lower-triangularity, graph-support, stability, diagonal-innovation, and positive-definiteness constraints in the formally specialized compatible-fiber definition. Hence
\[
\mathcal S(B,Q;G)=
\begin{cases}
\{(B,Q)\},&\text{if all the displayed fine-grid compatibility constraints hold},\\
\varnothing,&\text{otherwise}.
\end{cases}
\]

It remains to verify the certificate output. Under \(m=1\), every enumerated diagonal-root equation is
\[
a_{ii}^{,1}=B_{ii},
\]
and therefore has the unique real solution \(a_{ii}=B_{ii}\). By the declared definition of the divided-difference polynomial,
\[
C_1(x,y)=\sum_{r=0}^{0}x^r y^{0-r}=1.
\]
A nondecreasing walk of length one has at most one strict increase, while \(R_{vu,1}\) contains only monomials indexed by walks having at least two strict increases. Thus every remainder used by the specialized triangular recursion is zero. At each allowed off-diagonal coordinate, the root equation consequently reduces to
\[
B_{ij}=A_{ij}C_1(A_{ii},A_{jj})+0=A_{ij}.
\]
There is therefore no singular divided-difference coordinate, and the graph-recursive part of Definition~\(\texttt{def:certificate}\) produces the unique transition candidate \(A=B\). Given that candidate, the coarse covariance equation is
\[
Q=\sum_{r=0}^{0}A^rD(A^r)^\top=D,
\]
so it produces the unique innovation candidate \(D=Q\). The certificate's exact compatibility test accepts this candidate exactly when it satisfies all conditions defining \(\mathcal S(B,Q;G)\). By the fiber characterization above, the certificate therefore returns incompatibility exactly when the fiber is empty and otherwise accepts the unique compatible pair \((B,Q)\).

In the latter case, any two compatible copies used in the certificate's transfer-inequality test both equal \((B,Q)\), so transfer-distinct witnesses cannot exist. Definition~\(\texttt{def:certificate}\) must therefore return the common transfer. By the declared definition of \(h_{uv}\), this transfer is
\[
h_{uv}(z)=\frac{A_{vu}z}{1-A_{vv}z}
          =\frac{B_{vu}z}{1-B_{vv}z}.
\]
Compatibility supplies \(|B_{vv}|<1\). Since \(z\in\mathbb T\) implies \(|z|=1\),
\[
|B_{vv}z|=|B_{vv}|<1,
\]
and hence \(1-B_{vv}z\ne0\). The displayed transfer is therefore well defined on all of \(\mathbb T\). Because \(A=B\), it is exactly the direct fine-grid transfer with no temporal aliasing. This proves the proposition.
